\chapter{Pianificazione}
\label{app:pianificazione}

Il progetto dura dodici settimane, dal 22/10/2026 al 13/01/2027. Ogni milestone è un controllo: se non è superato, non si va avanti. Il dettaglio è in \texttt{report/fasi-e-milestone.md}.

\begin{table}[htbp]
\centering
\caption{Le milestone.}
\label{tab:milestone}
\begin{tabularx}{\linewidth}{@{}llL@{}}
\toprule
\textbf{Milestone} & \textbf{Entro il} & \textbf{Cosa deve essere pronto} \\
\midrule
M1 & 04/11/2026 & Decisioni e previsioni scritte \\
M2 & 11/11/2026 & L'applicazione funziona (V1--V3) \\
M3 & 25/11/2026 & L'applicazione gira su Kubernetes; primo guasto misurato \\
M3b & 02/12/2026 & Prova su AWS: cluster creato e distrutto, costi noti \\
M4 & 09/12/2026 & Il banco di prova dà misure affidabili \\
M5 & 16/12/2026 & Tutti gli esperimenti locali fatti \\
M6 & 30/12/2026 & Esperimenti su AWS fatti, tutto spento \\
M7 & 13/01/2027 & Consegna \\
\bottomrule
\end{tabularx}
\end{table}

\paragraph{Impegno stimato.} Circa 306 ore in due, cioè circa 13 ore a settimana a testa.

\paragraph{Se siamo in ritardo} rinunciamo, in quest'ordine, a: spegnimento di una macchina (F6), V5, rallentamento del database (F5), parte degli esperimenti su AWS, una delle due modalità di V4. Non rinunciamo mai a V1--V3, ai guasti principali e alle previsioni.
